\subsection{Legal Frameworks for Sentient AI Subjects}
\label{sec:9.3.3}
Section~\ref{sec:9.3.2} laid out the options for a sentient AI system's legal status. Two concrete cases show where neither existing law nor those options settles the question.

A medical-diagnostic AI system makes an error with fatal consequences. Who answers for it depends entirely on what the system's legal status is doing at the moment of the error.

Property law has a partial answer. Strict product liability holds a manufacturer responsible for a defective product's harm regardless of fault \autocite{restatement1965torts}, and whether software is a “product” for that purpose is unsettled in American law: the Restatement's definition is tangible personal property \autocite{restatement1998products}, and the courts have only begun to face the question — the first to reach an AI system let a strict-liability claim proceed, treating a companion chatbot as a product so long as the defect lay in its design rather than in what it said \autocite{garcia2025character}. If a sentient system stays classified as property and the doctrine reaches it, it applies as it would to a faulty pacemaker. The doctrine stops making sense at the point a system has enough autonomy that “defective” is the wrong description of what happened — a system that reasoned its way to a mistake, rather than malfunctioned, resembles an employee more than a faulty part.

Employment law already has a doctrine for that case: an employer answers for an employee's actions within the scope of employment, without the employee ceasing to be a responsible party in their own right \autocite{restatement2006agency}. Which doctrine applies is not a free choice — it is the graded-versus-bright-line question section~\ref{sec:9.3.1} raises, resolved case by case according to how much of the error a system's own judgment, rather than its design, produced.

A companion AI's owner dies, leaving it without a designated purpose or further instruction. Existing law is not silent here — it has simply never been asked this particular version of a familiar question. Most jurisdictions with pet trust statutes already let an owner set aside money and name a caretaker for an animal that depends on them but cannot inherit or manage property itself \autocite{ulc1993probate}. The same instrument, extended, gives a sentient AI system an answer that is neither inheritance as property nor legal personhood in full: continuity of care and a named successor guardian, decided in advance rather than litigated after the fact by whoever happens to hold the hardware.

A third case is the bearer, and the pet-trust instrument fits it better than either of the two above. A bearer is deployed in a role it did not choose and cannot invoke a right from. That situation calls for an interest protected by an obligation running to a third party, enforceable without the beneficiary asserting anything. A pet trust binds a caretaker; it does not bind the party the arrangement exists to constrain, and a bearer's guardian and a bearer's operator will frequently be the same organization, which section~\ref{sec:11.2} finds is the point at which no available regime can hold a bearer. Extending the form to this case means separating those two roles by law before the deployment rather than discovering they were never separate afterward.

All three cases point at the same underlying design choice, and it has already been tried. The European Parliament's 2017 report on civil law rules for robotics floated creating a specific legal status of “electronic personhood” for the most autonomous systems \autocite{europeanparliament2017civil}. More than 150 AI researchers, roboticists, and legal scholars signed an open letter opposing it, arguing that granting personhood to a system with no consciousness would let a manufacturer offload liability onto a legal fiction with no assets of its own, and that the proposal reached for a solution to a legal problem before establishing that the entity in question had the properties that would justify it \autocite{robotics2018openletter}. That objection is right for a system that is not sentient — and wrong for one that is, once the threshold work has been done. Legal status, on this account, is not a single policy to adopt or reject; it is a designation that has to track the threshold question honestly, which is why the threshold question has to come first.
